\documentclass[10pt]{article}
\providecommand{\FIGDIR}{fig}
\newcommand{\DIGESTDATE}{22 October 2026}
\input{preamble}

\begin{document}

% ==================================================================== MASTHEAD
\thispagestyle{empty}
\begin{tikzpicture}[remember picture, overlay]
  \fill[Deep] (current page.north west) rectangle
      ($(current page.north east)+(0,-67mm)$);
  \fill[Amber] ($(current page.north west)+(0,-67mm)$) rectangle
      ($(current page.north east)+(0,-68.4mm)$);
  \foreach \x/\y/\r in {22/12/0.9, 41/26/1.5, 63/9/0.7, 88/31/1.1, 112/17/1.9,
                        138/29/0.8, 159/13/1.3, 178/54/1.0, 196/21/1.6,
                        30/44/0.6, 74/57/1.2, 125/48/0.9, 168/6/0.8, 191/59/0.7,
                        52/18/1.0, 100/60/0.6, 148/20/1.1, 66/49/1.4, 205/45/0.8}
    \fill[Sky, opacity=0.30] ($(current page.north west)+(\x mm,-\y mm)$) circle (\r mm);
\end{tikzpicture}

\vspace*{4mm}
{\sffamily\fontsize{7.6}{9}\selectfont\color{Sky}%
 AUTOMATED DAILY LITERATURE DIGEST\;\textbullet\;PhD RESEARCH WATCH\par}
\vspace{2.2mm}
{\sffamily\bfseries\fontsize{27}{29}\selectfont\color{white}%
 Particulate Matter\par}
{\sffamily\fontsize{27}{29}\selectfont\color{Sky}%
 Monitoring \& Health Effects\par}

\vspace{5.0mm}
{\sffamily\fontsize{8.4}{10}\selectfont\color{white}%
 Issue 2026-10-22\;\;\textcolor{Amber}{\textbar}\;\;Entry window 22 October 2026 (full day)\;\;%
 \textcolor{Amber}{\textbar}\;\;17 records screened in scope\par}
\vspace{18mm}

\noindent
\begin{minipage}[t]{0.190\linewidth}\metric{17}{RECORDS IN SCOPE\\(38 EXCLUDED)}{Deep}\end{minipage}\hfill
\begin{minipage}[t]{0.190\linewidth}\metric{8}{OF 10 SUBTOPIC\\BANDS POPULATED}{Teal}\end{minipage}\hfill
\begin{minipage}[t]{0.190\linewidth}\metric{6}{EFFECT ESTIMATES\\WITH A CI}{Sky}\end{minipage}\hfill
\begin{minipage}[t]{0.190\linewidth}\metric{4}{TIER-A\\RECORDS}{Amber}\end{minipage}\hfill
\begin{minipage}[t]{0.190\linewidth}\metric{1}{RECORDS CARRIED ON\\METADATA ALONE}{Clay}\end{minipage}

\vspace{6mm}

\section*{\sffamily\bfseries\color{Deep}Signal of the day}
\vspace{-1.5mm}{\color{Amber}\rule{\linewidth}{1.1pt}}\vspace{2.5mm}

\begin{keybox}
{\sffamily\bfseries\small\color{Deep}Denmark, 53,901 people: the ADHD signal sits in the second year of
life, it follows elemental carbon rather than PM$_{2.5}$ mass, and it is larger in girls.}\\[1.4mm]
\footnotesize
A national case-cohort (21,860 diagnosed cases, 32,041 subcohort; births 1981--2008) with daily address-level PM$_{2.5}$
and EC from pre-conception to age 12 (DEHM/UBM/AirGIS). Each period was adjusted for the others. Second year of life,
EC $>$1.13 vs $<$0.49\,$\mu$g\,m$^{-3}$: HR \textbf{1.24} (1.12--1.36), unchanged after PM$_{2.5}$ adjustment;
girls \textbf{1.44} (1.21--1.72), boys 1.18 (1.00--1.40). Models adjust for ADHD polygenic scores.
\textbf{Where it is weak:} exposures at one address in consecutive years are close to collinear, so a ``window''
found by mutual adjustment is sensitive to differential error between periods. ADHD diagnosis in girls carries
referral bias that tracks urbanicity.
\textbf{Why it matters here:} the combustion tracer, not total mass, carries the effect. A mass-only optical network
cannot see this contrast; an absorption (BC) channel can.\par
\end{keybox}

\vspace{3mm}

\section*{\sffamily\bfseries\color{Deep}What changed today}
\vspace{-1.5mm}{\color{Amber}\rule{\linewidth}{1.1pt}}\vspace{2.5mm}

\footnotesize
\begin{itemize}

\item \textbf{HAPIN null extends to age 5.} LPG stove plus fuel cut infant PM$_{2.5}$ from 113 to 37\,$\mu$g\,m$^{-3}$,
but HAZ $\beta$ $-$0.06 ($-$0.16 to 0.03) across 2,169 children and no stunting effect; post-trial exposures
converged ($\approx$99\,$\mu$g\,m$^{-3}$ both arms). Preprint.

\item \textbf{Blocked cross-validation as interpolation.} GEOS-CF PM$_{2.5}$ emulator, Arabian Peninsula: a Random
Forest on latitude/longitude alone (R$^2$ 0.91) beat the 13-predictor model (0.82--0.84). Skill vanished with a
300\,km buffer and on an unseen sector ($r$ 0.11).

\item \textbf{CO$_2$-normalised traffic emission factors.} Helsinki street canyon: PM and NO$_x$ fleet factors well
above post-2014 Euro limits; traffic is a \emph{sink} for O$_3$ and sub-2\,nm ions (negative ``impact factors'').

\item \textbf{Intermittent point source.} Gas-fired peaker plant, Massachusetts: respiratory ED/UC visits OR 1.044
(1.001--1.090) per unit of a unitless trajectory index; no cardiovascular signal.

\item \textbf{Implausible per-microgram hazard.} Beijing elderly cohort ($n$\,=\,3,472): all-cause HR 1.053 per
1\,$\mu$g\,m$^{-3}$, an order of magnitude above pooled long-term estimates --- read as urban/rural confounding.

\item \textbf{Europe PMC lag.} Europe PMC returned a single 22 October record (a medRxiv preprint); PubMed (21) and the
Crossref ISSN sweep (37) carried the issue.

\end{itemize}

\vspace{2mm}

\begin{keybox}
{\sffamily\bfseries\footnotesize\color{Deep}Window continuity}\\[1.2mm]
{\fontsize{7.2}{8.8}\selectfont
The 21 October issue closed with \texttt{last\_entry\_date} at 21 October. This issue collects
\textbf{22 October in full}, continuous with it. The scheduled 04:04 UTC run fired at 16:06 UTC on 23 October
(11:06 CDT), before the 22:00 cut-off, so 23 October is \emph{not} opened and is left to the next run at or after
22:00 CDT. The W43 weekly (18--24 October) is due Saturday 24 October.\par}
\end{keybox}

\vspace{3mm}
\par\vskip 0pt plus 18\baselineskip\penalty-200\vskip 0pt plus -18\baselineskip
\section*{\sffamily\bfseries\color{Deep}Corpus analytics}
\vspace{-1.5mm}{\color{Amber}\rule{\linewidth}{1.1pt}}\vspace{2.5mm}

\noindent\begin{minipage}{\linewidth}\includegraphics[width=\linewidth]{\FIGDIR/f1_subtopics.png}
\figcap{Subtopic assignment of the 17 in-scope records across eight bands. Exposure assessment leads with
five (emission factors, emulation, soluble PM$_{10}$, a rainfall event, a microbiome method); respiratory three;
occupational/indoor and other clinical endpoints are empty.}\end{minipage}

\vspace{2mm}
\noindent\begin{minipage}{\linewidth}\includegraphics[width=\linewidth]{\FIGDIR/f2_design_endpoint.png}
\figcap{Left --- architecture: four reviews, three measurement campaigns, two cohorts, two
modelling, two toxicology, one trial, one case-crossover, one ecological; \textbf{metadata-only 1 (6\,\%)}, a series
low. Right --- endpoints run from ADHD and child growth to mortality, ED visits, CRS, ILD and in-vitro cytotoxicity.}\end{minipage}

\vspace{2mm}
\noindent\begin{minipage}{\linewidth}\includegraphics[width=\linewidth]{\FIGDIR/f3_metric_geography.png}
\figcap{Left --- particle metric: PM$_{2.5}$-only 9, unspeciated/emissions 4, joint
PM$_{2.5}$+PM$_{10}$ 2, PM$_{10}$ 1, nanoscale 1. Right --- geography: Europe 3 (Denmark, Finland, Spain) ties
China 3; global/multi-region 5 (mostly reviews); North America 2; one each for MENA, South Asia and Latin America.}\end{minipage}

\vspace{2mm}
\noindent\begin{minipage}{\linewidth}\includegraphics[width=\linewidth]{\FIGDIR/f6_lifecourse.png}
\figcap{Life-course windows: childhood carries two records (ADHD second-year window; HAPIN growth to age
5), with HAPIN also in utero; three working-age records; one older-adult cohort (Beijing); no adolescent record.}\end{minipage}

\vspace{2mm}
\noindent\begin{minipage}{\linewidth}\includegraphics[width=\linewidth]{\FIGDIR/f4_heatmap.png}
\figcap{Subtopic $\times$ study architecture. Exposure assessment is spread over measurement, modelling and
methods cells; the health bands sit in cohorts, the trial and reviews. Sensing has only the metadata-only
fluorescence record today.}\end{minipage}

\vspace{2mm}

\noindent\begin{minipage}{\linewidth}\includegraphics[width=\linewidth]{\FIGDIR/f5_forest.png}
\figcap{Ratio scale, five estimates from three studies: second-year EC and ADHD (all, female) in the Danish
case-cohort; PM$_{2.5}$ per 1\,$\mu$g\,m$^{-3}$ and all-cause/CVD death in the Beijing elderly cohort; peaker-plant
index and respiratory ED/UC visits (OR per unitless index unit). Exposure contrasts differ completely. Not poolable.}\end{minipage}

\vspace{3mm}

\noindent\begin{minipage}{\linewidth}\includegraphics[width=\linewidth]{\FIGDIR/f5_forest_2.png}
\figcap{Difference scale: HAPIN intention-to-treat effect of the LPG intervention on repeated height-for-age Z
score, ages 2--5 (n\,=\,2,169). The null is an ITT contrast for a partial exposure reduction.}\end{minipage}

\vspace{4mm}

\par\vskip 0pt plus 24\baselineskip\penalty-200\vskip 0pt plus -24\baselineskip
\section*{\sffamily\bfseries\color{Deep}Paper-level digest}
\vspace{-1.5mm}{\color{Amber}\rule{\linewidth}{1.1pt}}\vspace{2.5mm}

{\fontsize{7.4}{9}\selectfont\color{Slate}Ordered by cluster. Each entry gives the
methodological core and the finding that matters, not an abstract paraphrase. Records
marked \textit{metadata only} carry no recoverable abstract and assert no finding.\par}

\band{Neuro / mental health\hfill 1 record}{Deep}

\paperentry{Kalkbrenner traffic air pollution ADHD Denmark critical windows 2026}
{Traffic-related air pollution links with attention deficit hyperactivity disorder in Denmark: Developmental periods of susceptibility and sex-specific associations}
{Environ Res}
{Danish singletons born 1981-2008: 21,860 ADHD cases and a 32,041-person random subcohort. Exposure in each developmental period was mutually adjusted for the other periods to separate windows. The second year of life dominated: EC >1.13 vs <0.49 $\mu$g\,m$^{-3}$ HR 1.24 (1.12-1.36), not attenuated by PM$_{2.5}$ adjustment; females HR 1.44 (1.21-1.72) vs males 1.18 (1.00-1.40). PM$_{2.5}$ gave the same window and sex pattern. Polygenic scores for ADHD were adjusted. Limitation: period-specific exposures at the same address are highly collinear, and a window identified by mutual adjustment is sensitive to differential measurement error between periods. Diagnosis-based ADHD in females is also subject to referral bias that varies with urbanicity. Relevance: medium - EC, not PM$_{2.5}$ mass, carries the signal, which argues for absorption (BC) channels in dense networks.}
{10.1016/j.envres.2026.125900}

\par\vskip 0pt plus 10\baselineskip\penalty-200\vskip 0pt plus -10\baselineskip
\band{Cardiovascular \& metabolic\hfill 2 records}{Teal}

\paperentry{Wan PM$_{2.5}$ mortality urban rural elderly Beijing cohort 2026}
{Association between fine particulate matter exposure and all-cause and cardiovascular mortality in urban and rural elderly population in Beijing}
{Zhonghua Liu Xing Bing Xue Za Zhi}
{Beijing Elderly Comprehensive Health Cohort, n=3,472. Per 1 $\mu$g\,m$^{-3}$ PM$_{2.5}$: all-cause mortality HR 1.053 (1.026-1.080), CVD mortality HR 1.046 (1.001-1.100). Outdoor physical activity modified the association (P-interaction 0.022): HR 1.087 per $\mu$g\,m$^{-3}$ without outdoor activity vs 1.026 with it. Limitation: a 5\% hazard increase per single $\mu$g\,m$^{-3}$ is roughly an order of magnitude above the pooled long-term estimates. In a city-scale cohort the PM$_{2.5}$ contrast is mostly urban vs rural, so residual confounding by healthcare access and SES is the likely driver. The physical-activity 'offset' is also confounded by baseline health. Relevance: low - single-city satellite exposure.}
{10.3760/cma.j.cn112338-20251204-00870}

\paperentry{Yadav PM pollution-accelerated cardiovascular aging PACA framework 2026}
{Particulate Matter Exposure and Accelerated Cardiovascular Aging: Mechanistic Evidence, Biomarker Integration, and the PACA Framework}
{Cardiovasc Toxicol}
{A conceptual review that proposes the 'Pollution-Accelerated Cardiovascular Aging' (PACA) framework to separate transient toxicity from a real shift in cardiovascular ageing trajectory. It collects established pathways (mitochondrial and redox dysfunction, epigenetic change, senescence/SASP, proteostasis, endothelial and autonomic dysfunction, the pulmonary-neural-cardiac axis) and adds a Temporal-Evidence Matrix that weighs persistence after exposure stops and change relative to age-referenced comparators. Limitation: hypothesis-generating with no new data or systematic search. Without validated ageing clocks for the vasculature the matrix is hard to operationalise. Relevance: low - longitudinal personal exposure data are the missing input the framework asks for.}
{10.1007/s12012-026-10209-5}

\par\vskip 0pt plus 10\baselineskip\penalty-200\vskip 0pt plus -10\baselineskip
\band{Respiratory \& allergic\hfill 3 records}{Sky}

\paperentry{Jung peaker power plant emissions ED urgent care visits Massachusetts 2026}
{Individual-level health effects of living near a gas-fired "Peaker" power plant}
{J Clim Change Health}
{Primary-care patients living within 30 km of a gas-fired peaker plant, 2017-2021; 2,310 ED/UC visits analysed with a case-crossover design. Exposure was a daily HyADS parcel-trajectory index (mean 1.4, SD 16, reflecting plant intermittency); models adjusted for heat index and non-peaker PM$_{2.5}$ (mean 8.5 $\mu$g\,m$^{-3}$). Respiratory visits OR 1.044 (1.001-1.090) per unit; non-white patients OR 1.347 (1.034-1.755). Cardiovascular, headache and pregnancy visits were null. Limitation: the exposure is unitless and ZIP-code assigned, so the OR has no concentration scale. Peaker 'on' days coincide with peak heat and grid stress. A lower CI of 1.001 across four outcome families invites multiplicity concerns. Relevance: high - intermittent point sources are exactly what fence-line LCS can time-resolve.}
{10.1016/j.joclim.2026.100762}

\paperentry{Argersinger air pollution chronic rhinosinusitis review 2026}
{The link between air pollution and chronic rhinosinusitis}
{Curr Opin Otolaryngol Head Neck Surg}
{A short review of recent work on PM$_{2.5}$ and chronic rhinosinusitis. PM$_{2.5}$ is associated with CRS prevalence and severity and with worse disease-specific and quality-of-life outcomes after endoscopic sinus surgery. In nasal epithelial cells from CRS patients it worsens remodelling and barrier dysfunction and reduces ciliogenesis. Traffic, industrial, wildfire and pollen-related PM contribute independently. Limitation: non-systematic, no pooled estimates. Most cited associations are cross-sectional or claims-based with modelled residential exposure, and CRS phenotype (with or without polyps) is rarely separated. Relevance: low.}
{10.1097/moo.0000000000001164}

\paperentry{Kim interstitial lung disease changing climate Personal View 2026}
{Interstitial lung diseases in a changing climate: emerging risks and opportunities to improve care}
{Lancet Respir Med}
{A Lancet Respiratory Medicine Personal View on climate-related exposures in interstitial lung disease. Air pollution is the established risk factor. Extreme heat, viral illness, damp housing, pollen, disasters and desertification are framed as emerging threats. Wildfire smoke, occupational ILD under climate change and disaster preparedness are named as the main gaps. Limitation: an opinion synthesis with no quantitative evidence grading. ILD cohorts are small and heterogeneous, so most climate-ILD links are extrapolated from COPD and asthma. Relevance: low-medium - wildfire-smoke alerting for ILD patients is a concrete use case for hyperlocal sensing.}
{10.1016/s2213-2600(26)00252-3}

\par\vskip 0pt plus 10\baselineskip\penalty-200\vskip 0pt plus -10\baselineskip
\band{Reproductive \& developmental\hfill 1 record}{Sage}

\paperentry{Kalam HAPIN LPG intervention child growth stunting ages 2-5 2026}
{Effects of a 500-day LPG cookstove and fuel intervention delivered during gestation and infancy on height-for-age (HAZ) and stunting at ages 2, 3, 4 and 5 years}
{medRxiv (preprint)}
{HAPIN randomised trial (\textasciitilde{}3,200 pregnant women; LPG stove plus free fuel to infant age 12 months), extended follow-up of 2,169 children. The intervention cut gestational PM$_{2.5}$ from 117.5 to 80.7 $\mu$g\,m$^{-3}$ and infant PM$_{2.5}$ from 113.2 to 36.9 $\mu$g\,m$^{-3}$, but post-trial exposures converged (99.9 vs 98.7). ITT effect on repeated HAZ: beta -0.06 (-0.16 to 0.03); no effect on stunting at any age and no effect modification. Limitation: even the intervention arm stayed far above the WHO guideline during gestation, so this is a null for a partial exposure contrast, not for clean air. The exposure-response analysis that would test dose is still pending. Preprint, not peer reviewed. Relevance: medium - the trial's personal monitoring shows how far stove switching moves exposure.}
{10.64898/2026.10.16.26365712}

\par\vskip 0pt plus 10\baselineskip\penalty-200\vskip 0pt plus -10\baselineskip
\band{Mechanistic toxicology\hfill 2 records}{Violet}

\paperentry{Amato mobile air-liquid interface toxicity Barcelona harbor urban 2026}
{Assessment of air pollution toxicity in urban environments in Barcelona, Spain, using a mobile air-liquid interface system}
{Environ Sci Pollut Res Int}
{A mobile air-liquid interface system exposed A549 cells directly to ambient air at a harbour site and an urban background site, alongside co-located particle and gas monitoring. Partial least squares explained up to 80\% of the variance in ALI toxicity from air-quality parameters. Gas-phase pollutants, ozone above all, contributed more than particles. Limitation: LDH release in a single tumour-derived cell line is a crude endpoint. With two sites and highly collinear predictors, PLS loadings cannot separate ozone from photochemical-age covariates. n and exposure durations are not given in the abstract. Relevance: medium - pairing in-situ bioassays with multipollutant monitors is the template for testing whether mass-based LCS metrics track toxicity.}
{10.1007/s11356-026-38306-6}

\paperentry{Jiang carbon nanoparticle NSCLC progression ARNT2 ARL14 2026}
{Carbon nanoparticle exposure promotes the progression of non-small cell lung cancer via the ARNT2/ARL14 axis: Findings from in vivo and in vitro models}
{Ecotoxicol Environ Saf}
{Mouse CNP-induced fibrosis model (200 ug/kg every other day, 90 days) identified ARNT2 by transcriptomics. In xenografts with continued CNP exposure, tumour volume at day 36 rose (A549 194 to 308 mm3; A427 245 to 417 mm3) along with Ki-67/PCNA and N-cadherin/vimentin. ARNT2 knockdown reversed these changes. In vitro, 2 ug/mL for 48 h increased proliferation, migration (up to 3.6-fold), invasion, stemness markers and IL-6 (up to 13.7-fold), and reduced apoptosis. Limitation: carbon black nanoparticles stand in for 'the shared core' of PM$_{2.5}$, coal dust and smoke, but lack their organic and metal fractions. Instillation doses and in-vitro concentrations far exceed lung deposition from ambient air. Group sizes are not reported. Relevance: low.}
{10.1016/j.ecoenv.2026.120951}

\par\vskip 0pt plus 10\baselineskip\penalty-200\vskip 0pt plus -10\baselineskip
\band{Exposure assessment \& modelling\hfill 5 records}{Clay}

\paperentry{Savolainen road traffic emission factors impact factors Helsinki street canyon 2026}
{From direct emission factors to inverse and indirect impact factors of road traffic: influence of emerging and unregulated pollutants}
{Atmos Chem Phys}
{Wintertime supersite measurements in a central Helsinki street canyon. Fleet emission factors were derived by normalising pollutant increments to CO$_2$ as a fuel-consumption tracer, across particle number and size, BC, PM$_{10}$, PM$_{2.5}$, organic and inorganic PM components, NOx and VOCs including PAHs. PM and NOx emission factors were well above the post-2014 Euro limits, so older and heavy-duty vehicles dominate. Traffic showed negative 'impact factors' for O$_3$ and sub-2 nm small ions and contributed terpenoid VOCs. Limitation: one canyon in one season; the CO$_2$-ratio method assumes co-emission and common dilution, which fails for secondary and non-exhaust species. Relevance: high - a CO$_2$-normalised increment is a method LCS nodes with NDIR CO$_2$ can apply directly.}
{10.5194/acp-26-14823-2026}

\paperentry{Alturki ensemble ML emulation GEOS-CF PM$_{2.5}$ Arabian Peninsula spatial transfer 2026}
{Interpretable Ensemble Machine Learning for Statistical Emulation of Summer PM$_{2.5}$ Change in GEOS-CF over the Arabian Peninsula: How Far Does Spatially Blocked Skill Transfer?}
{Atmosphere}
{Emulated the 2020-2025 change in summer PM$_{2.5}$ over 3,933 GEOS-CF cells from 13 satellite and reanalysis predictors. Tree ensembles reached spatially blocked CV R2 0.82-0.84 (linear 0.55), but a Random Forest on latitude and longitude alone scored higher (R2 0.91). The predictor model lost all skill once cells within 300 km were excluded, failed on an unseen eastern sector (r 0.11) and did not transfer to other summers. Plain interpolation filled gaps better than any emulator. Split-conformal 90\% intervals covered 82-98\% of test cells, but only 60\% in the east. Limitation: the target is model output, not observations. Relevance: high - a clean demonstration that blocked-CV R2 in PM mapping can be pure interpolation; the same diagnostics belong in sensor-fusion papers.}
{10.3390/atmos17101035}

\paperentry{Bloise water-soluble carbon nitrogen PM$_{10}$ Chennai street vertical 2026}
{Seasonal, Spatial, and Near-Surface Vertical Variability of Water-Soluble Carbon and Nitrogen in PM$_{10}$ in Chennai, India}
{Atmosphere}
{Three Chennai street environments across three campaigns, sampled at nominal heights of 3, 16 and 37 m: 76 PM$_{10}$ samples, 75 WSOC/WSTN and 74 WSIN determinations. Means: PM$_{10}$ 65.3 $\mu$g\,m$^{-3}$, WSOC 2.53 ug C/m3, WSTN 1.28 and WSIN 1.26 ug N/m3. WSOC fell from 4.01 to 0.62 ug C/m3 between campaigns while WSIN rose from 1.12 to 1.86. Site contrasts were strong for PM$_{10}$, and vertical differences were not significant. Limitation: about 25 samples per campaign split across sites and heights gives little power for the vertical null. Campaigns confound season with meteorology. Relevance: medium - no vertical gradient between 3 and 37 m supports rooftop siting of street-canyon LCS for PM$_{10}$, though the null is underpowered.}
{10.3390/atmos17101033}

\paperentry{Magalhaes extreme rainfall landslides PM$_{10}$ PM$_{2.5}$ Sao Sebastiao 2026}
{Can extreme rainfall associated with landslides affect air quality? A case study in a coastal city (Sao Sebastiao, Sao Paulo State, Brazil)}
{Environ Monit Assess}
{Hourly PM and meteorology from December 2021 to February 2025, with k-means meteorological clusters and openair polar and time-variation analysis around the February 2023 extreme rainfall and landslides. Rain first scavenged PM. The post-landslide period then showed PM$_{10}$ up 124\% and PM$_{2.5}$ up 93.3\%. Two comparison rain events showed similar scavenging but different recovery, so the authors do not attribute the rise to landslides alone. Limitation: a single station and one event; clean-up traffic, resuspension of dried debris and meteorology are not separated, and there is no composition data. Relevance: medium - post-disaster resuspension is a case for rapidly deployable sensor nodes.}
{10.1007/s10661-026-16071-w}

\paperentry{Kim SpaMixed spatial mixed models microbiome PM$_{2.5}$ 2026}
{Spatial mixed models for assessing environmental exposure effects on the microbiome}
{Genome Med}
{SpaMixed is a spatial mixed model for microbiome count data that handles both spatial dependence across regions (conditional autoregressive priors) and ecological correlation among taxa. In simulations it achieved high feature-selection power with low false-positive rates and lower estimation error than existing methods. Applied to two PM$_{2.5}$ studies, it found 'biologically plausible' PM$_{2.5}$-associated taxa. Limitation: the abstract gives no sample sizes, regions or effect sizes for the applications. CAR priors absorb the spatial variance that carries the PM$_{2.5}$ contrast, so the method trades confounding for attenuation. Relevance: low - statistical method, no sensing content.}
{10.1186/s13073-026-01770-3}

\par\vskip 0pt plus 10\baselineskip\penalty-200\vskip 0pt plus -10\baselineskip
\band{Sensing, forecasting \& instrumentation\hfill 1 record}{Amber}

\paperentry{Cover solvatochromic fluorophores in situ aerosol fluorescence 2026}
{Evaluation of solvatochromic fluorophores for in situ fluorescence characterization of aerosols}
{Aerosol Sci Technol}
{\textsc{metadata only.} Title and authors only (Cover, Gibbons, Boadu, Ohno et al.). No abstract was retrievable from Crossref or Semantic Scholar at harvest. The title describes evaluating solvatochromic dyes, whose emission shifts with local polarity, as in-situ probes of aerosol particle properties. Analytes, particle types and performance metrics cannot be stated. Limitation: unverifiable beyond the title. Solvatochromic readouts also confound polarity with viscosity and water content, so any phase-state inference needs independent calibration. Relevance: low-medium - a laboratory composition probe, not a field-deployable sensor channel.}
{10.1080/02786826.2026.2746515}

\par\vskip 0pt plus 10\baselineskip\penalty-200\vskip 0pt plus -10\baselineskip
\band{Burden, policy \& mitigation\hfill 2 records}{Coral}

\paperentry{He ozone-dominant air pollution risk ageing inequality China counties 2026}
{A nationwide effect-based assessment of ozone-dominant air pollution risk patterns and ageing-related inequality in China}
{Ecotoxicol Environ Saf}
{Interpretable machine learning was fit to cancer data from 813 Chinese counties (2018-2019) to get non-linear long-term exposure-response shapes for cold-season PM$_{2.5}$ and warm-season O$_3$. Model-based thresholds then anchored exceedance and excess-hazard indices projected over 2017-2023. PM$_{2.5}$ pressure fell while O$_3$ pressure rose, with hotspots in Sichuan-Chongqing and the Yangtze River Delta that coincide with ageing. Limitation: two years of county cancer data cannot identify a dose-response threshold. ML-derived 'thresholds' in an ecological panel are model artefacts confounded by region, and the abstract gives no effect sizes or CIs. Relevance: low - county-level modelled exposure.}
{10.1016/j.ecoenv.2026.120945}

\paperentry{Scott equine air quality index Delphi management 2026}
{Experts recommend management strategies for horses during poor air quality}
{J Am Vet Med Assoc}
{A three-round modified Delphi with 15 equine veterinarians, 11 of whom completed all rounds. The panel adapted the EPA AQI public messaging to equine health risk and set management, training and competition recommendations for each AQI category. Experts agreed that healthy horses as well as compromised ones can have short- and long-term effects. Limitation: 11 completers, and the recommendations rest on translational extrapolation from human data, as the authors acknowledge. There are no equine exposure-response data. Relevance: medium - venue-level LCS readings at stables and competition grounds are what would operationalise AQI-tiered decisions.}
{10.2460/javma.26.05.0436}

\clearpage
\section*{\sffamily\bfseries\color{Deep}Corpus register}
\vspace{-1.5mm}{\color{Amber}\rule{\linewidth}{1.1pt}}\vspace{2.5mm}

\input{table_rows}

\vspace{2mm}

\begin{keybox}
{\sffamily\bfseries\footnotesize\color{Deep}Provenance and method}\\[1.2mm]
{\fontsize{7.2}{8.8}\selectfont
Entry window \textbf{22 October 2026}, single day (Thursday). Harvester legs: \textbf{PubMed} health
\textbf{21}, sensing \textbf{1}; \textbf{Europe PMC} \textbf{1} (index lag); \textbf{Crossref by ISSN} \textbf{37};
\textbf{arXiv} 150 kB, no in-window entry; \textbf{OpenAlex} HTTP 429. 58 unique, 8 already carried, 50 fresh;
PubMed connector (EDAT 22 Oct) 25 PMIDs, 5 appended, \textbf{55 screened}. Abstracts for 9 PubMed records recovered
by efetch and for 6 more via Crossref/Semantic Scholar. \textbf{Consensus} calibration sweep: 10 hits, 0 new.
\textbf{17 admitted}, \textbf{38 rejected} with logged reasons. \textbf{DOI gate} 0 fail / 1 warn (CMA DOI live at
doi.org, not in Crossref). \textbf{Trial watch} 22--23 October: 2 registrations updated (APCIQ-BP HEPA purifier RCT
NCT05903950, AIRWEIGHS NCT02763917), both already tracked. \textbf{1 carried on metadata alone}. \textbf{Scholar
Gateway} not called. Abstracts remain the copyright of their publishers.\par}
\end{keybox}

\end{document}
